\section{One compatible arithmetic estimate}
\label{rh-budget}

The preceding results meet at the discrepancy $\overline E$ in
\eqref{rh-balance}. It must be estimated by the geometric and
integral maps that produce that identity. The same section weights,
field embeddings, periods, ranks and finite levels must occur in
every comparison. This section gives the final sufficient estimate
without requiring each intermediate error to satisfy a separate
worst-case bound.

\subsection{Only weighted lower comparisons are needed}

For each outer index $u$, let a finite nonempty set of maps have
positive rank weights $r_{u,i}$, and write $\langle\cdot\rangle_u$
for their rank-weighted mean. Let $\nu_{u,i}$ be quotient-source
slopes, $\mu_{u,i}$ initial section slopes, $h_{u,i}$ map-height
bounds, $e_{u,i}$ exponent terms, $D_{u,i}$ divisor terms and
$l_{u,i}$ quotient losses. Write $c_u$ for the common target slope.
The following lemma is a real-variable rearrangement, so all these
quantities must be in the same normalization.

\begin{lemma}[The combined lower comparison]\label{rh-lower-comparison}
Suppose
\[
 \langle\nu-\mu+e+l\rangle_u\ge\Gamma_u+F_u,
 \quad c_u+\langle h\rangle_u-F_u\le\langle D\rangle_u,
\]
\[
 s_u=\langle\nu-h\rangle_u,\qquad s_u-c_u\le P_u.
\]
If $\overline E=\operatorname{avg}_u(A_u-\mathcal M_u)$, then
\begin{equation}\label{rh-nonframe}
 \overline E\le\mathcal B_0+\operatorname{avg}_uP_u,\qquad
 \mathcal B_0=\operatorname{avg}_u
 \{A_u-\Gamma_u+\langle D+e+l-\mu\rangle_u-\mathcal M_u\}.
\end{equation}
The outer average may use any fixed positive probability weights.
\end{lemma}

\begin{proof}
The first two assumptions imply
$s_u\ge\Gamma_u+c_u-\langle D+e+l-\mu\rangle_u$.
Combine with $s_u-c_u\le P_u$, add $A_u-\mathcal M_u$, and average.
\end{proof}

For the arithmetic application, the outer weights are those of the
same translation average as \eqref{rh-balance}. The quantity $P_u$
may be the minimum of the valid slope bounds in
Section~\ref{rh-energy}. A convention using squared-norm slopes
causes no ambiguity: divide slopes and heights by two, apply the
ordinary slope theorem, and multiply back. The displayed penalties
are homogeneous of degree one in the slope spread and height gap.

The nonlinear penalties should be averaged with their actual
parameters. In particular, the retained source and target spaces,
the mass $M_u$, the source spread $U_u$ and target gap $H_u$ need not
be replaced by unrelated extrema. Vanishing trace errors can make
$P_u$ vanish, but do not remove $\mathcal B_0$. Its counting,
divisibility, quotient and section-degree terms still need estimates.

\subsection{A concrete scalar envelope}

The following calculation makes the final exponent loss explicit.
It is conditional on the arithmetic envelope itself. Retain the
notation of Section~\ref{rh-transfer}, put $\tau=t_0/n$, and use
$K>0$ for an auxiliary rank parameter, distinct from $\kappa$.

\begin{proposition}[Four normalized terms]\label{rh-four-terms}
Suppose $x\ge1$, $\epsilon>0$, $C\ge0$, $0\le a\le2xb$, and
\[
 \overline E\le x^{C\epsilon}
 \left[t_0a(1+m^2/K)
       +m\left(2nmbx^\epsilon+\frac{dKbx^{2\epsilon}}n\right)\right].
\]
Then
\begin{equation}\label{rh-four-terms-bound}
 \frac{\overline E}{\kappa}\le x^{C\epsilon}
 \left[\frac{4x}{dm}+\frac{4xm}{dK}
       +\frac{4mx^\epsilon}{d\tau}
       +\frac{2Kx^{2\epsilon}}{\tau n^2}\right].
\end{equation}
If, in addition,
\[
 n\ge x^{2+8\epsilon},\quad
 x^{1/2-\epsilon/16}\le m\le2x^{1/2-\epsilon/16},
\]
\[
 x^{9/2+10\epsilon}\le K\le2x^{9/2+10\epsilon},\quad
 \tau\ge\tfrac12x^{-\epsilon/8},\quad d\ge x^{-\epsilon},
\]
then
\begin{equation}\label{rh-thirty-six}
              \frac{\overline E}{\kappa}
                   \le36x^{1/2+(C+33/16)\epsilon}.
\end{equation}
\end{proposition}

\begin{proof}
Divide by $\kappa=t_0mbd/2$ and use $t_0=\tau n$.
The last term's density $d$ cancels. Before the common factor
$x^{C\epsilon}$, the four terms are bounded respectively by
\[
 4x^{1/2+17\epsilon/16},\quad
 8x^{-3-145\epsilon/16},\quad
 16x^{1/2+33\epsilon/16},\quad
 8x^{1/2-31\epsilon/8}.
\]
For $x\ge1$, their exponents are at most $1/2+33\epsilon/16$;
their coefficients sum to $36$.
\end{proof}

The factor $36$ is harmless at the analytic endpoint. What must
tend to zero is the exponent loss, including $C\epsilon$.
If $C$ varies with the accuracy, that product must still tend to
zero. A uniform finite constant during each large-scale limit is
enough; it need not be absorbed into a further power of $x$.

The numerical ranks are constrained by the construction. For example,
the first and third terms of \eqref{rh-four-terms-bound} balance at
$m=\sqrt{\tau x^{1-\epsilon}}$. If
$\tau\le2x^{-\epsilon/8}$, this is at most
$\sqrt2x^{1/2-9\epsilon/16}$, eventually below the displayed lower
window for $m$. A formal minimization of the envelope alone therefore
does not provide a permissible replacement family.

\subsection{The final acceptance theorem}

\begin{theorem}[A single signed budget suffices for RH]
\label{rh-final}
For each $j$, fix $\beta_j>1$, $h_j\ge0$ and
$1/2<\theta_j<1$, with $\theta_j\to1/2$. Let
$\mathcal Y_j\subset(1,\infty)$ be unbounded.
Suppose the same arithmetic data give the hypotheses of
Corollary~\ref{rh-direct-budget} and Lemma~\ref{rh-lower-comparison}
for almost every $x\in[Y,e^{h_j}Y^{\beta_j}]$, for each
$Y\in\mathcal Y_j$. At each such $x$ define
\begin{equation}\label{rh-final-budget}
 \begin{split}
 \mathcal B_j(x)={}&
 \frac{\mathcal B_0(x)+\operatorname{avg}_uP_u(x)}{\kappa(x)}
 +\frac1{t_0(x)}\sum_tC_{q,\ell}(a_t/b)\\
 &+E_{\rm thick}(x)+E_{\rm thin}(x)+(1+\log(x+1))^2.
 \end{split}
\end{equation}
Assume $\mathcal B_j$ is measurable and locally integrable. If
\[
 \liminf_{\substack{Y\to\infty\\Y\in\mathcal Y_j}}
 \frac{\log\left(1+\displaystyle\int_Y^{Y^{\beta_j}}
    (\mathcal A_{h_j}\mathcal B_j(x))_+
                       x^{-\theta_j}\,\frac{dx}{x}\right)}{\log Y}=0
\]
for every $j$, then RH holds. The signed $-J/\kappa$ may replace
$E_{\rm thin}$ throughout when retained in the direct comparison.
\end{theorem}

\begin{proof}
The two comparisons give $G\le\mathcal B_j$ on the expanded
intervals. The positive averaging operator and the positive-part
function preserve that inequality. Apply Corollary~\ref{rh-averaged}
along the same unbounded sets of starting points.
\end{proof}

For a simpler sufficient version, take $h_j=0$ and establish a
pointwise bound $\mathcal B_j(x)\le M_jx^{\theta_j}$ on the whole
windows. Proposition~\ref{rh-four-terms} is one possible supplier
for the discrepancy part, once its input envelope is proved.
A compatible quadratic estimate
\[
 \frac1{t_0}\sum_tC_{q,\ell}(a_t/b)
                           \le M_{c,j}x^{1/2+\delta_{c,j}}
\]
with $\delta_{c,j}\to0$ then leaves the explicit transfer and
centering errors. Absolute estimates for the three quadratic sums
are sufficient, but a bound for their signed combination is enough.
All sample arguments must lie in the interval on which the chosen
quadratic estimates hold, including the shifts of at most one.

The uniform analytic inputs can also be kept smaller. Suppose a single
$\vartheta_0<1$ is known such that every Dirichlet $L$-function is
zero-free in $\Re s>\vartheta_0$, apart from the principal pole.
For every real zero $\beta\in(0,1)$ of a primitive real character of
conductor $q\ge3$, this immediately gives
\[
 (1-\beta)\log q\ge(1-\vartheta_0)\log3>0.
\]
Thus this particular uniform exclusion of zeros close to one requires
no separate no-Siegel assumption on a route already using the common
zero-free half-plane. It does not assert the absence of all real
zeros in the strip, or replace the separate geometric estimates.

The arithmetic characters, groups and lattices may vary with $x$.
The analytic target throughout this RH manuscript remains the single
fixed function $G$. By contrast, the accuracy-dependent interval
exponent, averaging width and analytic exponent must remain fixed
as the admissible windows tend to infinity.

\subsection{A finite version of the final construction}

The arithmetic choices can instead be made only at integers. For each
fixed $j$, take arbitrarily large integers $N$ and supply the comparisons
in Theorem~\ref{rh-final} at every integer
$N\le k\le\lceil e^{h_j}N^{\beta_j}\rceil$.
Let $B_{j,N}(k)$ be the resulting right side of
\eqref{rh-final-budget}. These tables may be chosen independently for
different intervals. If
\[
 \liminf_{N\to\infty}
 \frac{\log\!\left(1+\sum_{n=N}^{\lfloor N^{\beta_j}\rfloor}
       \frac{(\mathsf A_{h_j}B_{j,N}(n))_+}{n^{1+\theta_j}}\right)}
      {\log N}=0
\]
along those chosen integers, then RH follows from
Lemma~\ref{shared-integer-table}. Indeed each table bounds $G$ at its
sampled integers, and the interpolation error tends to zero.
This formulation requires finitely many compatible constructions per
interval and no measurable choice of arithmetic data. It retains all
the arithmetic estimates needed at those integers.

\subsection{The exact mathematical frontier}

Three construction tasks remain. The first is a uniform lower bound
for the actual normalized test intersection in
Corollary~\ref{geo-tower-block-bound}, together with that corollary's
other comparison hypotheses, the prescribed period and finite-valuation
identifications, and the quantitative finite section-degree estimates.
Section~\ref{rh-geometry} then converts that input into the needed
degree contribution using one mean volume and short sections.

The second is a positive normalization and a common-target identity
for the actual integral theta maps, followed by bounds for their
two positive spectral traces and their arithmetic heights. The matrix
identities in Section~\ref{rh-energy} specify sufficient quantities;
they do not establish their size for those maps.

The third is to assemble the remaining counting, divisibility,
quotient, auxiliary-character and cohomological estimates with those
same fields, levels and scales. Uniform quasi-RH and no-Siegel
inputs can assist the auxiliary analytic and arithmetic estimates
\cite{rh-qrh,rh-siegel}; a fixed zero-free boundary to the right of
one half does not itself give the vanishing exponent loss in
Theorem~\ref{rh-final}. The full signed budget, rather than a list
of independently small errors, is the remaining sufficient target.
